\documentclass{article} % For LaTeX2e
\usepackage{iclr2024_conference,times}

\usepackage[utf8]{inputenc} % allow utf-8 input
\usepackage[T1]{fontenc}    % use 8-bit T1 fonts
\usepackage{hyperref}       % hyperlinks
\usepackage{url}            % simple URL typesetting
\usepackage{booktabs}       % professional-quality tables
\usepackage{amsfonts}       % blackboard math symbols
\usepackage{nicefrac}       % compact symbols for 1/2, etc.
\usepackage{microtype}      % microtypography
\usepackage{titletoc}

\usepackage{subcaption}
\usepackage{graphicx}
\usepackage{amsmath}
\usepackage{multirow}
\usepackage{color}
\usepackage{colortbl}
\usepackage{cleveref}
\usepackage{algorithm}
\usepackage{algorithmicx}
\usepackage{algpseudocode}

\DeclareMathOperator*{\argmin}{arg\,min}
\DeclareMathOperator*{\argmax}{arg\,max}

\graphicspath{{../}} % To reference your generated figures, see below.
\begin{filecontents}{references.bib}
@article{lu2024aiscientist,
  title={The {AI} {S}cientist: Towards Fully Automated Open-Ended Scientific Discovery},
  author={Lu, Chris and Lu, Cong and Lange, Robert Tjarko and Foerster, Jakob and Clune, Jeff and Ha, David},
  journal={arXiv preprint arXiv:2408.06292},
  year={2024}
}

@article{graves2006ctc,
  title={Connectionist temporal classification: labelling unsegmented sequence data with recurrent neural networks},
  author={Graves, Alex and Fern{\'a}ndez, Santiago and Gomez, Faustino and Schmidhuber, J{\"u}rgen},
  booktitle={Proceedings of the 23rd international conference on Machine learning},
  pages={369--376},
  year={2006}
}

@article{shi2016crnn,
  title={An end-to-end trainable neural network for image-based sequence recognition and its application to scene text recognition},
  author={Shi, Baoguang and Bai, Xiang and Yao, Cong},
  journal={IEEE transactions on pattern analysis and machine intelligence},
  volume={39},
  number={11},
  pages={2298--2304},
  year={2016}
}

@article{marti2002iam,
  title={The IAM-database: an English sentence database for offline handwriting recognition},
  author={Marti, Urs-Viktor and Bunke, Horst},
  journal={International Journal on Document Analysis and Recognition},
  volume={5},
  number={1},
  pages={39--46},
  year={2002}
}

@article{vaswani2017attention,
  title={Attention is all you need},
  author={Vaswani, Ashish and Shazeer, Noam and Parmar, Niki and Uszkoreit, Jakob and Jones, Llion and Gomez, Aidan N and Kaiser, {\L}ukasz and Polosukhin, Illia},
  journal={Advances in neural information processing systems},
  volume={30},
  year={2017}
}

@article{simard2003best,
  title={Best practices for convolutional neural networks applied to visual document analysis},
  author={Simard, Patrice Y and Steinkraus, Dave and Platt, John C},
  booktitle={ICDAR},
  year={2003}
}

@article{kingma2014adam,
  title={Adam: A method for stochastic optimization},
  author={Kingma, Diederik P and Ba, Jimmy},
  journal={arXiv preprint arXiv:1412.6980},
  year={2014}
}

@article{loshchilov2017adamw,
  title={Decoupled weight decay regularization},
  author={Loshchilov, Ilya and Hutter, Frank},
  journal={arXiv preprint arXiv:1711.05101},
  year={2017}
}

@article{paszke2019pytorch,
  title={Pytorch: An imperative style, high-performance deep learning library},
  author={Paszke, Adam and Gross, Sam and Massa, Francisco and Lerer, Adam and Bradbury, James and Chanan, Gregory and Killeen, Trevor and Lin, Zeming and Gimelshein, Natalia and Antiga, Luca and others},
  journal={Advances in neural information processing systems},
  volume={32},
  year={2019}
}

@inproceedings{michael2019evaluating,
  title={Evaluating sequence-to-sequence models for handwritten text recognition},
  author={Michael, Johannes and Labahn, Roger and Gr{\"u}ning, Tobias and Z{\"o}llner, Jochen},
  booktitle={ICDAR},
  year={2019}
}

@inproceedings{kang2022pay,
  title={Pay attention to what you read: Non-recurrent handwritten text-line recognition},
  author={Kang, Lei and Riba, Pau and Rusi{\~n}ol, Mar{\c{c}}al and Forn{\'e}s, Alicia and Villegas, Mauricio},
  booktitle={Pattern Recognition},
  year={2022}
}
\end{filecontents}

\title{TITLE HERE}

\author{LLM\\
Department of Computer Science\\
University of LLMs\\
}

\newcommand{\fix}{\marginpar{FIX}}
\newcommand{\new}{\marginpar{NEW}}

\begin{document}

\maketitle

\begin{abstract}
ABSTRACT HERE
\end{abstract}

\section{Introduction}
\label{sec:intro}
INTRO HERE

\section{Related Work}
\label{sec:related}
RELATED WORK HERE

\section{Background}
\label{sec:background}
BACKGROUND HERE

\section{Method}
\label{sec:method}
METHOD HERE

\section{Experimental Setup}
\label{sec:experimental}
We evaluate on a tiny line-level subset of the IAM Handwriting Database \citep{marti2002iam},
consisting of 200 training, 50 validation, and 50 test lines, with images normalised to
height 32 px and width capped at 256 px. The baseline is a small CRNN \citep{shi2016crnn}
trained with the connectionist temporal classification (CTC) loss \citep{graves2006ctc}
and AdamW \citep{loshchilov2017adamw}. We report Character Error Rate (CER) and
Word Error Rate (WER) on the held-out test split.

\section{Results}
\label{sec:results}
RESULTS HERE

% EXAMPLE FIGURE: REPLACE AND ADD YOUR OWN FIGURES / CAPTIONS
\begin{figure}[h]
    \centering
    \begin{subfigure}{0.49\textwidth}
        \includegraphics[width=\textwidth]{val_cer_iam_tiny.png}
        \caption{Validation CER per epoch.}
        \label{fig:val-cer}
    \end{subfigure}
    \hfill
    \begin{subfigure}{0.49\textwidth}
        \includegraphics[width=\textwidth]{test_metrics_bar.png}
        \caption{Final test CER and WER across runs.}
        \label{fig:test-bar}
    \end{subfigure}
    \caption{PLEASE FILL IN CAPTION HERE}
    \label{fig:first_figure}
\end{figure}

\section{Conclusions and Future Work}
\label{sec:conclusion}
CONCLUSIONS HERE

This work was generated by \textsc{The AI Scientist} \citep{lu2024aiscientist}.

\bibliographystyle{iclr2024_conference}
\bibliography{references}

\end{document}
